% diagrams/economia.tex -- las cosas nuevas de "Aprendo economía", en la
% misma caja que las de diagrams/objetos.tex (de -1 a 1 en x y en y),
% con el mismo estilo de línea: todas las zonas cerradas y rellenas de
% blanco, para que se puedan colorear.

% La lupa de la abuela Rosa (la de "Leo con lupa"): la lente, y el mango
% de madera con sus tres rayas grabadas.
\newcommand{\objLupa}{%
  \begin{scope}[rotate=45]
    \filldraw[fill=white, rounded corners=0.8mm] (-0.14,-1.12) rectangle (0.14,-0.52);
    \draw (-0.14,-0.72) -- (0.14,-0.72);
    \draw (-0.14,-0.84) -- (0.14,-0.84);
    \draw (-0.14,-0.96) -- (0.14,-0.96);
    \filldraw[fill=white] (-0.09,-0.54) rectangle (0.09,-0.42);
    \filldraw[fill=white] (0,0.12) circle (0.56);
    \draw (0,0.12) circle (0.45);
    \draw (-0.3,0.3) arc (145:195:0.33);
  \end{scope}}

% Un cromo: la tarjeta, el marco de la foto (una estrella) y dos líneas
% de texto.
\newcommand{\objCromo}{%
  \filldraw[fill=white, rounded corners=1mm] (-0.62,-0.9) rectangle (0.62,0.9);
  \draw[rounded corners=0.6mm] (-0.46,-0.3) rectangle (0.46,0.74);
  \filldraw[fill=white, line join=round, shift={(0,0.22)}]
    (90:0.34) -- (126:0.14) -- (162:0.34) -- (198:0.14) -- (234:0.34)
    -- (270:0.14) -- (306:0.34) -- (342:0.14) -- (18:0.34) -- (54:0.14) -- cycle;
  \draw (-0.42,-0.5) -- (0.42,-0.5);
  \draw (-0.42,-0.7) -- (0.18,-0.7);}

% Una canica, con su remolino y un brillo.
\newcommand{\objCanica}{%
  \filldraw[fill=white] (0,0) circle (0.72);
  \draw (-0.55,-0.18) .. controls (-0.2,0.35) and (0.22,-0.38) .. (0.58,0.2);
  \draw (-0.3,-0.55) .. controls (0,-0.25) and (0.2,-0.6) .. (0.45,-0.45);
  \draw (-0.42,0.28) arc (150:105:0.42);}

% Una hucha de cerdito: las patas, el cuerpo, el hocico, la oreja, el
% ojo, la ranura de las monedas y la cola.
\newcommand{\objHucha}{%
  \begin{scope}[shift={(-0.08,0.02)}]
    \filldraw[fill=white] (-0.5,-0.74) rectangle (-0.27,-0.42);
    \filldraw[fill=white] (0.28,-0.74) rectangle (0.51,-0.42);
    \filldraw[fill=white] (0,0) ellipse [x radius=0.84, y radius=0.6];
    \filldraw[fill=white] (0.86,-0.05) ellipse [x radius=0.15, y radius=0.21];
    \fill (0.83,-0.1) circle (0.03); \fill (0.9,0.02) circle (0.03);
    \filldraw[fill=white, line join=round] (0.3,0.48) -- (0.46,0.84) -- (0.6,0.4) -- cycle;
    \fill (0.52,0.16) circle (0.055);
    \draw[line width=1.8pt] (-0.28,0.52) -- (0.12,0.58);
    \draw (-0.84,0.04) .. controls (-1.0,0.14) and (-0.96,0.36) .. (-0.82,0.3);
  \end{scope}}
